

\chapter{Clustern mit Machine Learning Algorithmen}
\label{chap:ml}

Zunächst betrachtet die Analyse Algorithmen aus dem Bereich des Machine Learning. Im Fokus stehen insbesondere die Cluster-Algorithmen DBSCAN, k-Means und Drain. Diese Algorithmen sind für produktive Anwendungen interessant, da freie und quelloffene Bibliotheken wie scikit-learn stabile und getestete Implementierungen bereitstellen.

Die Clusteranalyse mittels Machine-Learning-Algorithmen unterteilt sich in zwei Bereiche: partitionierende Verfahren wie K-Means und dichtebasierte Verfahren wie DBSCAN. Der folgende Abschnitt erläutert, warum sich partitionierende Verfahren für die kontinuierliche Log-Gruppierung weniger eignen.

\section{Numerische Textrepräsentation für maschinelle Verarbeitung}
Für die nachfolgende Verarbeitung ist es zunächst erforderlich, den Text in eine numerische Darstellung zu überführen. Hierbei stehen verschiedene Verfahren zur Auswahl, die sich hinsichtlich ihrer Komplexität und der Art der Informationserfassung unterscheiden.

\subsection{Bag-of-Words: Vorteile und Grenzen dimensionaler Vektorräume}

Ein grundlegender Ansatz ist das Bag-of-Words-Modell.\footnote{Der Artikel von IBM klärt über das Bag-of-Words-Modell auf: \url{https://www.ibm.com/de-de/think/topics/bag-of-words} (Zugriff: \today)} . Es basiert auf der Konstruktion eines Vektorraums, in dem jedes eindeutige Wort, das im Textkorpus vorkommt, eine eigene Dimension bildet. Umfasst das Vokabular beispielsweise n unterschiedliche Wörter, so weist der resultierende Vektorraum ebenfalls n Dimensionen auf. Innerhalb dieses mehrdimensionalen Raums wird jedes Textdokument als einzelner Datenpunkt verortet. Die Position des Punktes entlang einer bestimmten Achse (Dimension) wird dabei durch die Häufigkeit des entsprechenden Wortes im Dokument bestimmt. Das bedeutet, dass Begriffe, die besonders oft auftreten, einen stärkeren Einfluss auf die Positionierung des Dokuments im Vektorraum haben. Dies führt insbesondere bei umfangreichen Textsammlungen zu einer sehr hohen Dimensionalität, die sowohl Chancen als auch Herausforderungen für die weiterführende Analyse mit sich bringt. 


Risiken sind dabei:
\begin{itemize}
    \item Überanpassung (Overfitting): Wenn ein Dokument aus vielen unterschiedlichen Wörtern besteht, werden die Vektorräume sehr groß (viele Dimensionen), das führt zu Performance einbüßen und hohen Speicherverbrauch. 
    \item Verlust von Zusammenhängen: Je höher die Dimensionen durch Einzelwörter, desto schwieriger wird es ohne Zusatztechniken, semantische Zusammenhänge (Wort-Korrelationen) zu erkennen.
\end{itemize}

Daraus ergeben sich folgende Chancen. 
\begin{itemize}
    \item Feingranulare Informationen: Durch die hohe Dimensionalität geht kein Wort verloren.
    \item Einfachheit \& Effektivität: Trotz der hohen Dimensionen ist das Modell durchaus effektiv für Aufgaben wie Dokumentenklassifizierung (zum Beispiel Spam-Filter)
\end{itemize}
Eine wesentliche Einschränkung dieses Modells liegt zudem in der Vernachlässigung der Syntax. Da das Modell die grammatikalische Struktur und die Reihenfolge der Wörter ignoriert, gehen im Zuge der Abstraktion teilweise semantische Zusammenhänge verloren. Dieser Punkt ist bei Lognachrichten jedoch weniger relevant, da sich darin oft keine oder nur wenige semantische Zusammenhänge ergeben. 

\subsection{TF-IDF: Kontextbezogene Gewichtung technischer Log-Terme}


Im Kontext klassischer Algorithmen hat sich zur Realisierung einer ressourcenschonenden und zugleich präzisen Gewichtung das TF-IDF-Maß etabliert. Diese statistische Methode quantifiziert die Relevanz eines einzelnen Wortes innerhalb eines Dokuments in Relation zur gesamten Dokumentenkollektion.Ein wesentlicher Vorteil dieses Verfahrens liegt in der Überwindung der Schwächen reiner Häufigkeitszählungen: Während allgemein häufige Begriffe in ihrer Gewichtung herabgestuft werden, erfahren seltene, aber inhaltlich signifikante Terme – wie beispielsweise spezifische Fehlercodes – eine stärkere Gewichtung.Basis dieser Berechnung ist die sogenannte „Term Frequency“, welche die Häufigkeit eines Wortes in einem betrachteten Dokument beschreibt. Formal lässt sich dies wie folgt darstellen \cite[Kap.~11]{jm3}:

\begin{equation}
  \text{tf}_{t,d} = 
  \begin{cases} 
    1 + \log_{10} \text{count}(t,d) & \text{if } \text{count}(t,d) > 0 \\
    0 & \text{otherwise}
  \end{cases}
\end{equation}

Die Logarithmusfunktion vermeidet eine übermäßige Gewichtung von Wörtern mit hoher Vorkommensfrequenz. Der zweite Teil des Tf-IDF heißt IDF (Inverse Document Frequency) \cite[Kap.~11]{jm3}:
\begin{equation}
  \text{idf}_t = \log_{10} \frac{N}{\text{df}_t}
\end{equation}
Die IDF reduziert die Gewichtung von Wörtern, die in vielen Dokumenten auftreten, und verstärkt seltene Wörter. Dies trägt zur differenzierten Gewichtung einzelner Wörter bei.Das ist besonders praktisch bei der Verarbeitung von Log-Nachrichten, da dort viele technische Wörter vorkommen, denen jedoch nicht automatisch überdurchschnittlich viel Bedeutung zugesprochen werden muss. Diese ausgewogene Verarbeitung von Text ist der Grund für die weitere Verwendung zum Vektorisieren und Verarbeiten der Log-Nachrichten mit Methoden des maschinellen Lernens.

\subsection{Semantische Textmodelle: Kontextuelle Vektorisierung durch Embeddings}
Ein zentraler Baustein der modernen Sprachverarbeitung sind Embeddings – Vektorrepräsentationen, die die Bedeutung von sprachlichen Einheiten direkt aus deren Verteilung in Textdaten erlernen. Dieser Ansatz, der als Vektorsemantik bezeichnet wird, basiert auf der Distributionshypothese \cite[Kap.~5]{jm3}, nach der Wörter oder Texte, die in ähnlichen Kontexten auftreten, auch eine ähnliche Bedeutung tragen. Technisch betrachtet wird dabei jeder Text als Punkt in einem mehrdimensionalen semantischen Raum abgebildet. Die Stärke dieses Verfahrens liegt darin, dass semantische Ähnlichkeit in mathematische Distanz übersetzt wird: Je näher zwei Vektoren in diesem Raum beieinander liegen, desto ähnlicher ist ihre inhaltliche Bedeutung.
Während klassische Ansätze oft statische Embeddings für einzelne Wörter generieren, ist es für die Verarbeitung ganzer Textpassagen notwendig, den Kontext dynamisch zu erfassen \cite[Kap.~5.5]{jm3}. Zur praktischen Vektorisierung der Eingabedaten wird daher auf das Framework \texttt{sentence-transformers (SBERT)\footnote{Webseite der Python Bibliotek: \url{https://www.sbert.net/} (Zugegriffen: \today)}} zurückgegriffen (vgl. sbert.net). Dieses Framework ermöglicht die Generierung von Satz-Embeddings, die speziell für semantische Ähnlichkeitsvergleiche optimiert sind.
Konkret kommt in dieser Arbeit das vortrainierte Modell \texttt{sentence-transformers/all-MiniLM-L6-v2} zum Einsatz. Die Wahl fiel auf dieses Modell, da es eine optimale Ausgewogenheit zwischen der Vektorlänge (hier 384 Dimensionen) und der Verarbeitungsgeschwindigkeit (Inferenzzeit) bietet. Dies gewährleistet eine ressourcenschonende Berechnung, die dennoch eine hohe Präzision bei der Abbildung semantischer Zusammenhänge aufweist. Das gewählte Modell wurde mit verschiedenen Datensätzen angelernt, darunter auch ein großer Anteil an Code von Programmiersprachen. Zu den angestrebten Verwendungen dieses Modells zählt auch das Clustering von Text. Dabei ist jedoch zu beachten, dass das Modell eine maximale Wortlänge von 256 Zeichen hat. Alles, was darüber hinausgeht, wird abgeschnitten.\footnote{\url{https://huggingface.co/sentence-transformers/all-MiniLM-L6-v2} erläutert alle Details zu praktischen Verwendung des Modells} In den späteren Experimenten und Tests zur Gruppierung von Lognachrichten werden die Embedding-Modelle nur für die Verarbeitung mit Vektordatenbanken (Kapitel~\ref{chap:vek}) genutzt, da die nachfolgenden Ansätze mit Machine-Learning-Algorithmen sehr leichtgängig sind und durch die aufwendigere Vorverarbeitung ausgebremst werden würden. 

\section{k-Means: Partitionierendes Clustering und dessen Einschränkungen}
\label{sec:kmeans}
Nachdem die Nutzung der passenden Textvektorisierung geklärt ist, werden im Folgenden die Algorithmen zum Gruppieren bzw. Clustern der Log-Nachrichten beschrieben.

Der K-Means-Clusteralgorithmus umfasst drei Schritte. Die Nutzung erfordert zunächst die Wahl eines Wertes für k.

\begin{enumerate}
    \item Der Algorithmus bestimmt zunächst K Masseschwerpunkte im Datensatz. Diese Schwerpunkte heißen Centroide $c(k)$ für $k = 1$ bis $K$.
    \item Im zweiten Schritt berechnet er für jeden Vektor $x^n$ die Distanz zu den $K$ Masseschwerpunkten. Die Distanzfunktion hängt von der Problemstellung ab. Oft nutzt man die euklidische Distanz \cite{Wilmott2020-rx}:
    \begin{equation}
      \text{Distanz}^{(n,k)} = \sqrt{ \sum_{m=1}^{M} \left( x_m^{(n)} - c_m^{(k)} \right)^2 } \quad \text{für } k=1 \text{ bis } K
    \end{equation}
    Jeder Datenpunkt erhält den nächstgelegenen Masseschwerpunkt als Zuordnung.
    \item Für jedes entstandene Cluster berechnet der Algorithmus den Mittelwert, wodurch die Centroide in die Nähe des Cluster-Mittelpunktes rücken. Die Schritte 2 und 3 wiederholen sich bis zur Konvergenz.
\end{enumerate}
Die Auswahl der Zahl $K$ ist essenziell für das Clustering.
Der manuelle Aufwand bei einem System, das periodisch neue Logs einliest, ist jedoch zu hoch. Selbst eine Automatisierung mittels der Elbow-Methode zur Bestimmung von K erforderte bei jedem Durchlauf eine Neuberechnung. Zudem ist die Anzahl der Loggruppen im Voraus unbekannt. Deshalb ist ein Cluster-Algorithmus erforderlich, der ohne feste Eingabeparameter wie K auskommt.

\section{DBSCAN: Flexibles, dichtebasiertes Clustering ohne vordefinierte Cluster-Anzahl}
\label{sec:dbscan}

Der DBSCAN-Algorithmus gehört zu den dichtebasierten Clusterverfahren. Anders als k-Means benötigt er keinen Parameter k, was eine flexiblere Log-Gruppierung ermöglicht. DBSCAN definiert einen Cluster als dichte Region aus Datenpunkten, getrennt durch Bereiche geringerer Dichte. Der Algorithmus unterteilt Datenpunkte in drei Kategorien: Kernpunkte, Grenzpunkte und Rauschpunkte. Eine dichte Region definiert sich als Radius $\varepsilon$ um einen Punkt $x  \in \mathbb{R}^d$ \cite{Mohammed_Wagner_2020}:
\begin{equation}
    N_{\epsilon}(\mathbf{x}) = B_{d}(\mathbf{x}, \epsilon) = \{ \mathbf{y} \mid \| \mathbf{x} - \mathbf{y} \| \leq \epsilon \}
\end{equation}
Dabei bezeichnet $\| \mathbf{x} - \mathbf{y} \|$ die euklidische Distanz, ersetzbar durch andere Distanzfunktionen. Der Parameter $\epsilon$ gibt den Radius an. Zusätzlich definiert der Wert $minpts$ einen Schwellwert: Eine Anzahl von Punkten innerhalb eines Radius $\epsilon$ gilt ab diesem Wert als relevant. Ein Datenpunkt $x \in D$ heißt Kernpunkt, wenn er mindestens $minpts$ Nachbarn im Radius $\epsilon$ besitzt. Hat ein Punkt weniger Nachbarn, gilt er als Grenzpunkt. Liegt ein Datenpunkt zwar in Reichweite, aber nicht im $\epsilon$-Radius eines Kernpunktes, ist er ein Rauschpunkt. DBSCAN arbeitet in folgenden Schritten \cite{geeksforgeeksDBSCANClustering}:
\begin{enumerate}
    \item Kernpunkte identifizieren: Für jeden Datenpunkt $x \in D$ zählt der Algorithmus die Punkte innerhalb des Radius $\epsilon$. Übersteigt die Anzahl $minpts$, gilt der Punkt als Kernpunkt.
    \item Jeder noch nicht zugeordnete Kernpunkt bildet ein neues Cluster. Anschließend sucht der Algorithmus rekursiv nach allen dichte-verbundenen Punkten innerhalb des Radius $\epsilon$ und ordnet diese dem Cluster zu. Zwei Punkte $x$ und $y$ sind dichte-verbunden, wenn eine Kette von Punkten innerhalb des Radius $\epsilon$ existiert, von denen mindestens einer ein Kernpunkt ist.
    \item Alle Punkte ohne Cluster-Zugehörigkeit gelten als Rauschpunkte.
\end{enumerate}

\section{Drain3: Effizientes On-Demand-Parsing mittels Baumstrukturen}
\label{sec:drain3}


Im folgenden Abschnitt geht es um den Drain Online Log Parser, welcher in den algorithmischen Vergleichen durch die Python-Bibliothek drain3 \footnote{GitHub Repository zur Python Implementierung von drain: \url{https://github.com/logpai/Drain3}} genutzt wird. Die Bibliothek basiert auf dem Drain-Algorithmus. Die „3” im Namen zeigt an, dass es sich um eine Weiterentwicklung der Drain-Bibliothek handelt, welche für Python 3 optimiert wurde. Die Vorgänger-Bibliothek war für Python 2 entwickelt worden. 
Drain ist in der Lage, Lognachrichten kontinuierlich zu verarbeiten und ihnen dabei ein Logcluster sowie ein Logtemplate zuzuordnen. Ein Online-Logparser wie drain3 kann Logs eines Systems „on-demand“ verarbeiten, er muss nicht vorher auf Daten angelernt werden. Dies geschieht durch einen internen Parse-Baum mit fester Tiefe (Fixed Depth Tree), welcher den Kern des Algorithmus bildet. Statt Lognachrichten erst über einen Zeitraum zu sammeln und nachträglich auszuwerten, leitet Drain jede neu eintreffende Nachricht direkt durch diese Baumstruktur, um sie effizient dem passenden Cluster zuzuordnen. Der Prozess lässt sich vereinfacht in vier Schritte unterteilen \cite[Abschn.~III]{drain}:

\begin{enumerate}
    \item Vorverarbeitung: Zunächst werden mittels einfacher regulärer Ausdrücke offensichtliche Variablen wie IP-Adressen oder Block-IDs aus der Lognachricht entfernt, um die Genauigkeit zu verbessern.
    \item Suche nach Log-Länge: Auf der ersten Ebene des Baums sortiert Drain die Nachricht basierend auf ihrer Länge (Anzahl der Tokens). Die Annahme hierbei ist, dass Logs, die zum gleichen Event gehören, in der Regel auch die gleiche Anzahl an Wörtern besitzen.
    \item Suche nach Anfangs-Tokens: Auf den tieferen Ebenen navigiert der Algorithmus anhand der ersten Wörter der Nachricht. Da Logzeilen oft mit konstanten Begriffen beginnen, führen diese Pfade effizient in die richtige Richtung. Um eine Explosion des Baumes zu verhindern, werden Tokens, die Ziffern enthalten, hierbei nicht als eigener Pfad angelegt, sondern generalisiert.
    \item Ähnlichkeitsprüfung: Am Ende eines Pfades (im Blattknoten) befindet sich eine Liste potenzieller Loggruppen. Hier vergleicht Drain die neue Nachricht mit den dort gespeicherten Events. Wird ein passendes Template gefunden, wird die Nachricht der Gruppe hinzugefügt und das Template gegebenenfalls aktualisiert (z. B. durch Einfügen von Wildcards). Findet sich kein Match, erstellt Drain an dieser Stelle eine neue Loggruppe.
\end{enumerate}

\subsection{Formale Ähnlichkeitsbewertung: Die simSeq-Funktion}

Nachdem Drain im vierten Schritt den passenden Blattknoten erreicht hat, muss der 
Algorithmus entscheiden, ob eine neu eintreffende Lognachricht einem bereits 
existierenden Logtemplate zugeordnet werden kann oder ob eine neue Loggruppe 
anzulegen ist. Diese Entscheidung basiert auf der Ähnlichkeitsfunktion 
$simSeq$, welche im ursprünglichen Drain-Paper formal definiert ist
\cite[Abschn.~III.E]{drain}:

\begin{equation}
    simSeq = \frac{\sum_{i=1}^{n} equ(seq_1(i), seq_2(i))}{n}
    \label{eq:simseq}
\end{equation}

Dabei repräsentieren $seq_1$ die neu eintreffende Lognachricht und $seq_2$ das 
gespeicherte Logtemplate einer Kandidatengruppe. Der Ausdruck $seq(i)$ bezeichnet 
das $i$-te Token der jeweiligen Sequenz, während $n$ die Gesamtlänge beider 
Sequenzen in Token angibt – da Drain auf der ersten Baumebene bereits nach 
Log-Länge filtert, ist diese für beide Sequenzen identisch. Die Hilfsfunktion 
$equ$ bewertet den paarweisen Tokenvergleich:

\begin{equation}
    equ(t_1, t_2) = \begin{cases} 1 & \text{if } t_1 = t_2 \\ 0 & \text{otherwise} \end{cases}
    \label{eq:equ}
\end{equation}

$simSeq$ ergibt sich damit als der Anteil übereinstimmender Tokens an der 
Gesamtlänge der Nachricht – ein Wert zwischen $0$ und $1$. Wird der berechnete 
Wert mit dem konfigurierbaren Ähnlichkeitsschwellenwert $st$ verglichen 
(\lstinline{drain_sim_th} in der drain3-Bibliothek), gilt: Ist 
$simSeq \geq st$, wird die Nachricht der Kandidatengruppe mit dem höchsten 
$simSeq$-Wert zugeordnet und das Template gegebenenfalls durch Wildcards 
aktualisiert. Unterschreitet der beste Kandidat den Schwellenwert, legt Drain 
an dieser Stelle eine neue Loggruppe an.

Durch diese Struktur mit fester Tiefe vermeidet Drain, dass der Suchbaum zu tief und unbalanciert wird, was die hohe Geschwindigkeit bei gleichzeitiger Präzision ermöglicht.

Für den praktischen Einsatz der Python-Bibliothek drain3 ist zudem ein Aspekt relevant, der über die theoretische Beschreibung im ursprünglichen Paper hinausgeht: die Persistenz des Log-Modells. Der Drain Algorithmus ist als reines In-Memory-Verfahren konzipiert, bei dem der Parse-Baum flüchtig im Arbeitsspeicher existiert und theoretisch bei jedem Systemstart neu aufgebaut werden müsste. Die drain3-Bibliothek erweitert diesen Ansatz jedoch um eine Speicherfunktion (\textit{Snapshotting}). Dies ermöglicht es, den kompletten internen Zustand des Baumes inklusive aller gelernten Templates und Cluster auf die Festplatte zu schreiben und bei einem Neustart wiederherzustellen. In der Praxis muss der Parser somit nicht jedes Mal bei null beginnen ("Cold Start"), sondern kann nahtlos auf das historische Wissen zugreifen und neue Lognachrichten sofort effizient in die bestehende Struktur integrieren.

\subsection{Persistenz in Drain3: Zustandserhaltung des Parse-Baums}

Die Persistenzmechanik innerhalb der \lstinline{TemplateMiner-Klasse} \ref{lst:template_miner_py} der Datei \\ \texttt{templateminer.py} gewährleistet die Dauerhaftigkeit des trainierten Modellzustands über Laufzeiten hinweg. Die Architektur sieht hierfür die Nutzung eines abstrakten PersistenceHandlers vor, der bei der Initialisierung des TemplateMiners optional übergeben wird. Ist ein solcher Handler definiert, initiiert das System beim Start automatisch den Versuch, einen bereits existierenden Zustand mittels der Methode \lstinline{load_state()} wiederherzustellen. Der eigentliche Speichervorgang wird durch die Methode \lstinline{save_state()} abgebildet. Dabei verwendet die Funktion \lstinline{save_state} das Modul \lstinline{jsonpickle}, das ganze Python-Objektgraphen in das JSON-Format überführen kann. \lstinline{jsonpickle} ist eine Serialisierungsbibliothek, die beliebige Python-Objekte in lesbares JSON umwandeln und auch wieder exakt in ihre ursprüngliche Objektinstanz zurückführen kann. Dies gelingt, indem das Modul automatisch Python-spezifische Metadaten wie Klassennamen, Variablen und deren Werte in das JSON einbettet. Dadurch können selbst komplexe, verschachtelte Strukturen wie der Drain-Parse-Baum ohne Informationsverlust gespeichert und beim Neustart nahtlos wiederhergestellt werden. Die Funktion \lstinline{load_state} liest den Zustand aus der abgespeicherten Datei ein und dekomprimiert den Inhalt, wenn vorher eine Komprimierung mit \lstinline{zlib} stattgefunden hat.  Anschließend wandelt \lstinline{JSONPickle.load()} den Zustand wieder in eine Instanz der Drain-Klasse zurück. Danach werden Felder wie \lstinline{id_to_cluster} gesetzt. Dabei handelt es sich um ein Schlüssel-Wert-Paar, welches die Zuordnung von der Cluster-ID und den Logcluster-Daten aus dem Template, der Cluster-Größe und der ID ermöglicht. Auch der \lstinline{cluster_count} wird wiederhergestellt, der anzeigt, wie viele Cluster schon gefunden wurden, sowie das \lstinline{root_node}, der Wurzelknoten des Baums zur Verarbeitung der Logs. 
Die beiden Funktionen \lstinline{save_state()} und \lstinline{load_state()} stellen sicher, dass der Zustand des Parsers über einen Neustart hinweg gespeichert und wiederhergestellt werden kann. Durch die Implementierung als abstrakte Klasse bietet die Bibliothek viel Flexibilität, wie der Zustand des Systems gespeichert werden soll. Die Bibliothek kommt mit vorgefertigten PersistenceHandlern für Redis und Kafka, zwei Vertreter performanter Schlüssel-Wert-Datenbanken. 

\section{Brain: Log Parsing mit bidirektionalem Parallelbaum}
\label{sec:brain}

Brain \cite{brain} ist ein stabiler und effizienter Log-Parser, der auf der 
Beobachtung basiert, dass das längste gemeinsame Muster innerhalb einer Gruppe 
ähnlicher Lognachrichten mit hoher Wahrscheinlichkeit Teil des zugehörigen 
Log-Templates ist. Dieser Ansatz kombiniert Techniken der häufigen 
Elementextraktion mit einem heuristischen Baumverfahren und adressiert dabei 
eine zentrale Schwäche von Drain: Da Drain seinen Parse-Baum anhand der ersten 
Tokens einer Nachricht navigiert, versagt er bei Logs, deren erstes Token eine 
Variable ist – ein in der Praxis häufig auftretendes Muster, etwa bei 
Build-Systemen wie Jenkins, wo der Projektname den Anfang einer Logzeile bildet. 
Brain umgeht dieses Problem grundlegend, indem es nicht auf die Position der 
Tokens, sondern auf deren Häufigkeitsmuster zurückgreift.

\subsection{Funktionsweise}

Der Parsing-Prozess von Brain gliedert sich in fünf Schritte \cite[Abschn.~III]{brain}:

\begin{enumerate}
    \item \textbf{Vorverarbeitung:} Zunächst werden die Lognachrichten anhand 
    konfigurierbarer Trennzeichen in einzelne Tokens zerlegt. Bekannte Variablen 
    wie IP-Adressen werden mittels regulärer Ausdrücke durch Wildcards 
    \texttt{<*>} ersetzt. Für jedes Token wird anschließend dessen 
    Auftretenshäufigkeit über alle Logs hinweg erfasst. Tokens mit gleicher 
    Häufigkeit in einer Lognachricht werden zu sogenannten \textit{Wortkombinationen} 
    zusammengefasst.

    \item \textbf{Initiale Gruppenbildung:} Jede Lognachricht wählt aus ihren 
    Wortkombinationen diejenige mit den meisten Tokens als \textit{Longest Common 
    Pattern} aus. Lognachrichten gleicher Länge mit identischem Longest Common 
    Pattern werden einer initialen Gruppe zugeordnet. Da konstante Wörter 
    häufiger und gleichmäßiger auftreten als variable Wörter, ist das Longest 
    Common Pattern mit hoher Wahrscheinlichkeit Teil des gesuchten Log-Templates.

    \item \textbf{Knotenaktualisierung in der Eltern-Richtung:} Für jede initiale 
    Gruppe erstellt Brain einen bidirektionalen Baum, der am Longest Common 
    Pattern wurzelt. In der Eltern-Richtung werden Spalten untersucht, deren 
    repräsentative Häufigkeit über der des Wurzelknotens liegt. Tauchen in einer 
    solchen Spalte mehrere verschiedene Wörter auf, wird sie als Variable 
    klassifiziert, da konstante Positionen stets dasselbe Wort erzeugen.

    \item \textbf{Knotenaktualisierung in der Kind-Richtung:} In der 
    Kind-Richtung werden Spalten mit niedrigerer repräsentativer Häufigkeit 
    als der Wurzelknoten analysiert. Die Spalten werden nach der Anzahl 
    unterschiedlicher Wörter sortiert. Liegt diese Anzahl unterhalb eines 
    konfigurierbaren Schwellwerts, werden die Tokens der Spalte als konstant 
    eingestuft und die initiale Gruppe entsprechend in Untergruppen aufgeteilt. 
    Dieser Prozess wird rekursiv fortgesetzt, bis alle konstanten Wörter 
    identifiziert sind.

    \item \textbf{Template-Generierung:} Abschließend kombiniert Brain das 
    Longest Common Pattern mit den als konstant klassifizierten Knoten des 
    Baums zum vollständigen Log-Template. Variable Positionen werden dabei durch 
    Wildcards \texttt{<*>} repräsentiert.
\end{enumerate}

Durch die bidirektionale Baumstruktur erreicht Brain eine zeitliche Komplexität 
von nahezu $\mathcal{O}(n)$, wobei $n$ die Anzahl der Lognachrichten bezeichnet. 
In Experimenten auf 16 Benchmark-Datensätzen verarbeitet Brain eine Million 
Lognachrichten in durchschnittlich 46 Sekunden, gegenüber 146 Sekunden bei 
Drain \cite{brain}. Die durchschnittliche Group Accuracy beträgt dabei 0,981 
gegenüber 0,865 bei Drain – eine Verbesserung um 11,6~\% \cite{brain}.